\section{Reaction--diffusion: why does interest spread?}
\label{sec:reaction-diffusion}

\subsection{Two reservoirs and one idea}

In chemistry, a substance diffuses from high concentration to low concentration. D-Tracker borrows
that picture for interest: if a keyword receives strong attention in location group $j'$, some of
that attention may flow to group $j$ where activity is lower. The general reaction--diffusion
equation cited by the paper is
\begin{equation}
  \frac{\partial u_i}{\partial t}
  =f(u_i)+\sum_j D\,(u_j-u_i),
  \label{eq:rd-background}
\end{equation}
where $u_i$ is the concentration at position $i$ and $D$ is diffusion strength. The term $f(u_i)$ is
the \textbf{reaction}: local change at a position. The term $\sum_j D(u_j-u_i)$ is
\textbf{diffusion}: it is non-zero only when positions differ, and its sign says whether flow enters
or leaves.

\begin{center}
\begin{tikzpicture}[
  >=Latex,
  place/.style={circle, draw=DTNavy, thick, minimum size=18mm, fill=DTMist,
                align=center, font=\sffamily\small},
  flow/.style={->, very thick, DTTeal},
  label/.style={font=\small, align=center}
]
  \node[place] (b) at (0,0) {group $j'$\\$w_{ij'}=4$};
  \node[place] (a) at (4.2,0) {group $j$\\$w_{ij}=1$};
  \draw[flow] (b) to[bend left=12] node[above,label] {inflow because $4>1$} (a);
  \draw[flow, DTClaret] (a) to[bend left=12] node[below,label] {outflow if $1>4$} (b);
\end{tikzpicture}
\end{center}

\subsection{D-Tracker's latent version}

D-Tracker does not impose the equation on every keyword--location pair. It imposes it on a
\emph{latent core} of $d_k\cdot d_l$ dynamics. With $i=1,\ldots,d_k$ indexing keyword groups and
$j=1,\ldots,d_l$ indexing location groups, the paper writes
\begin{equation}
  \frac{dw_{ij}}{dt}
  =a_{ij}w_{ij}
   +\sum_{j'=1}^{d_l} d_{ijj'}\,(w_{ij'}-w_{ij}),
  \qquad 1\leq i\leq d_k,\;1\leq j\leq d_l.
  \label{eq:rd-latent}
\end{equation}
Two parameters need careful reading:
\begin{itemize}
  \item $a_{ij}$ is the \textbf{growth rate}. If $a_{ij}>0$, dynamics $w_{ij}$ tends to rise; if
        $a_{ij}<0$, it tends to fall. This is the ``reaction''.
  \item $d_{ijj'}$ is the \textbf{diffusion strength} from location group $j'$ to group $j$. The
        paper describes $\mathcal{D}\in\mathbb{R}_{+}^{d_k\times d_l\times d_l}$ as non-negative.
        The actual direction also depends on $(w_{ij'}-w_{ij})$: if $w_{ij'}>w_{ij}$, group $j$
        receives; if it is lower, group $j$ loses.
\end{itemize}

Here $w_{ij}$ is the activity of the dynamics for keyword group $i$ and location group $j$. Dynamics
from different keyword groups do not diffuse directly in this sum; they interact through the core
and the multilinear projection in Section~\ref{sec:model}. The full system has $d_k\cdot d_l$
equations with parameters $\mathbf{A}\in\mathbb{R}^{d_k\times d_l}$ and
$\mathcal{D}\in\mathbb{R}_{+}^{d_k\times d_l\times d_l}$.

\begin{keypoint}
$w>0$ makes ``activity'' readable; $a$ says rise or fall; $(w_{j'}-w_j)$ says flow direction; $d$
says how strong the flow is. This is why the model can be interpretable rather than a black-box
forecaster.
\end{keypoint}

\subsection{A concrete one-step calculation}

Suppose a dynamics in group $j$ has $w_j=1$, another dynamics in the same keyword group but location
group $j'$ has $w_{j'}=4$, $a_j=0.1$, and $d_{jj'}=0.5$. Ignoring other groups, the derivative is
\[
  \frac{dw_j}{dt}=0.1\cdot1+0.5\cdot(4-1)=0.1+1.5=1.6.
\]
With an Euler step $\Delta t=1$, the next value is approximately $1+1.6=2.6$. The dynamics is rising
because it has both growth and inflow from a higher group. If $w_{j'}=0$ instead, the diffusion term
becomes $0.5(0-1)=-0.5$: the same strength $d$ but the opposite direction.

\begin{plainwords}
Do not memorize the equation. Read three sentences: ``it grows or shrinks on its own'', ``it is
pulled toward another group'', and ``the larger the gap, the stronger the pull''. All three sentences live
in the two terms of Equation~\eqref{eq:rd-latent}.
\end{plainwords}

\subsection{Connecting back to observations}

The core $w_{ij}$ is latent. To connect it to an observed series $X_{uv}$ for real keyword $u$ and
real location $v$, the paper differentiates the projection and obtains
\begin{equation}
  \frac{dX_{uv}}{dt}
  =\sum_{i=1}^{d_k}\sum_{j=1}^{d_l}
   W^{(\mathrm{key})}_{iu}W^{(\mathrm{loc})}_{jv}
   \underbrace{\left(
     a_{ij}w_{ij}+\sum_{j'=1}^{d_l}d_{ijj'}(w_{ij'}-w_{ij})
   \right)}_{\text{latent reaction--diffusion system}}.
  \label{eq:observed-derivative}
\end{equation}
Read from right to left: the latent system creates $d_k d_l$ dynamics; the non-negative weights
$W^{(\mathrm{key})}_{iu}W^{(\mathrm{loc})}_{jv}$ decide which dynamics contribute to pair $(u,v)$.
Because the weights are non-negative, the projection cannot flip a sign: a rising dynamics still
contributes positively to its series. This is the bridge from the ``small core'' to the ``real series''.

\subsection{Interpretation is not causation}

The reaction--diffusion assumption says that a level gap can create flow. The model learns $a,d$ and
group factors to fit data; it does \emph{not} prove that a rise in the United States caused a rise in
Vietnam. A third factor (a global launch, news, social media) can create the same delayed pattern.
Read the result as a \textbf{structured diffusion hypothesis}, to be checked with domain knowledge and
additional experiments.

\begin{commonmistake}
``$d_{ijj'}>0$ means $j'$ always sends to $j$'' is incomplete. Flow also depends on
$w_{ij'}-w_{ij}$; if $w_{ij'}<w_{ij}$, the term is negative. The paper calls $d$ a strength, not a
fixed direction. Likewise, ``$a>0$'' only describes local tendency, not guaranteed growth forever.
\end{commonmistake}

\begin{selfcheck}
\begin{enumerate}
  \item With $w_j=2,w_{j'}=5,a_j=0,d_{jj'}=0.2$, what is the derivative?
  \item If $a_{ij}<0$ and the diffusion inflow is positive, must the dynamics fall? Why or why not?
  \item Describe a situation with a delay between two places that is not causal diffusion.
\end{enumerate}
\end{selfcheck}
